% Generated draft: requires mathematical and source review.
\documentclass[11pt]{article}
\usepackage[margin=25mm]{geometry}
\usepackage{fontspec}
\setmainfont{Latin Modern Roman}
\usepackage{amsmath,amssymb,mathtools,unicode-math}
\setmathfont{Latin Modern Math}
\usepackage{enumitem,needspace,xurl,hyperref}
\setlength{\parindent}{0pt}
\setlength{\parskip}{5pt}
\setlength{\emergencystretch}{3em}
\newcommand{\N}{\mathbb{N}}
\newcommand{\Z}{\mathbb{Z}}
\newcommand{\Q}{\mathbb{Q}}
\newcommand{\R}{\mathbb{R}}
\newcommand{\C}{\mathbb{C}}
\newcommand{\F}{\mathbb{F}}
\newcommand{\A}{\mathbb{A}}
\newcommand{\PP}{\mathbb{P}}
\newcommand{\E}{\mathbb{E}}
\newcommand{\eps}{\varepsilon}
\newcommand{\dd}{\,\mathrm{d}}
\begin{document}
\section*{UnsolvedMath ID: 701536}
\textbf{Source problem code:} catalogue entry
\subsection{Definitions and mathematical statement}
Regard a finitely generated group \(G\) as a discrete group. Let \(\widehat G\) denote its unitary dual: the set of unitary equivalence classes of irreducible unitary representations, equipped with the Fell topology. Write \(1_G\) for the trivial representation and set
\[
\widehat G_{\mathrm{fin}}=\{[\pi]\in\widehat G: \pi(G)\text{ is finite}\}.
\]
Following the supplied source, \(G\) has property \((\tau)\) if \([1_G]\) is isolated in \(\widehat G_{\mathrm{fin}}\) with its subspace topology; equivalently, there is an open subset \(U\subseteq\widehat G\) such that
\[
U\cap\widehat G_{\mathrm{fin}}=\{[1_G]\}.
\]
A finite simple group means a nontrivial finite group whose only normal subgroups are the identity subgroup and the whole group; the source imposes no nonabelian restriction.

\textbf{Question (A. Jaikin-Zapirain, 2014).} Does there exist a finitely generated group \(G\) with property \((\tau)\) such that
\[
\text{for every finite simple group }F,\quad
\text{there exists a surjective homomorphism }\varphi_F:G\twoheadrightarrow F?
\]
The same group \(G\) must work for every such \(F\).

\textbf{Supplied motivation and qualifications.} The record attributes to E. Breuillard, B. Green, M. Kassabov, A. Lubotzky, N. Nikolov, and T. Tao a theorem asserting that there exist \(\epsilon>0\) and \(k\in\N\) such that every finite simple group \(F\) has a generating set \(S\subseteq F\) with \(|S|\leq k\) for which the Cayley graph of \(F\) with respect to \(S\) is an \(\epsilon\)-expander. The constants are uniform over \(F\). The supplied record does not specify the expansion normalization or the convention for symmetrizing \(S\); this assertion is retained as contextual motivation.

The record also attributes to M. Ershov, A. Jaikin-Zapirain, and M. Kassabov, \emph{Mem. Amer. Math. Soc.}, \textbf{1186} (2017), the existence of a group with Kazhdan's property \((T)\) that maps surjectively onto every simple group of Lie type of rank at least \(2\). Here property \((T)\) means that the trivial representation is isolated in the full unitary dual. This contextual assertion is preserved as supplied: its wording does not explicitly specify finiteness of the Lie-type groups or the rank convention. These contextual assertions are supplied literature reports, not independently verified conclusions. The registry labels the main question \emph{open}; that label is also supplied rather than independently verified.
\subsection{Short English statement}
Does one finitely generated group with property \((\tau)\) have every finite simple group as a quotient? Here property \((\tau)\) requires the trivial representation to be isolated among irreducible unitary representations with finite image. The question is attributed to A. Jaikin-Zapirain (2014); its supplied status is \emph{open}.
\subsection{Sources}
\begin{itemize}
\item Local UnsolvedMath registry, permanent ID \(701536\), exact source code \texttt{catalogue entry}, title \emph{Kourovka Notebook Problem 18.108}. Supplied upstream revision: \texttt{d371ee917ac9647fbb06cd6fc529deec6f3dda82}. The registry supplies the status \emph{open}.
\item \emph{The Kourovka Notebook (Unsolved Problems in Group Theory)}, original Problem 18.108, 18th issue (2014); the record also identifies the 21st issue edition (2026), arXiv identifier \texttt{1401.0300v48}. The supplied collection description names E. I. Khukhro and V. D. Mazurov as editors. Proposed by A. Jaikin-Zapirain in 2014. This is the supplied primary-source attribution; no source was independently consulted.
\item E. Breuillard, B. Green, M. Kassabov, A. Lubotzky, N. Nikolov, and T. Tao: authors credited in the supplied record for the uniform-expander theorem used as motivation. No title, publication details, or link is supplied for this reference.
\item M. Ershov, A. Jaikin-Zapirain, and M. Kassabov, \emph{Mem. Amer. Math. Soc.}, \textbf{1186} (2017): supplied reference for the property \((T)\) quotient assertion. No title or link is supplied. The bibliographic designation \(1186\) is retained without assigning it a volume or issue interpretation.
\end{itemize}
\end{document}
